\section{Introduction}\label{sec:intro}

The study of pseudoscalar meson structure could shed light into the longstanding quantum-chromodynamics (QCD) mystery of the origin of mass.
In the Standard Model, the Higgs boson provides most of the masses of all noncomposite fermions.
However, the masses of the valence quarks of mesons contribute only a small fraction of the total meson mass.
The remainder is thought to arise from the dynamics of the strong interaction.
Studying kaon and pion structure can greatly aid our understanding of the emergence of hadronic mass.
One commonly studied type of structure is the parton distribution function (PDF), which gives the probability to find quarks and gluons inside a hadron as a function of their momentum fraction $x$ of the mother hadron.
Since the gluon field is responsible for almost all of the pion's mass, studying the gluon PDFs of mesons would be very valuable in understanding the structure of mesons.
There have been a number of attempts to extract pion gluon PDFs~\cite{Barry:2018ort,Cao:2021aci,Novikov:2020snp,Chang:2020rdy}, but not much progress has been made on the kaon PDFs.
The future experimental studies, for example, U.S.-based Electron-Ion Collider (EIC)~\cite{Accardi:2012qut,AbdulKhalek:2021gbh}, the Electron-Ion Collider in China (EicC)~\cite{Anderle:2021wcy},
the Drell-Yan and $J/\psi$-production experiments  COMPASS++/AMBER~\cite{Denisov:2018unj} in Europe  will aim at improving our knowledge of both the pion and kaon gluon and quark PDFs.
We refer readers to recent reviews~\cite{Aguilar:2019teb,Roberts:2021nhw,Arrington:2021biu} for more details on the current and future prospects for studies of pion and kaon structure.


Lattice QCD provides a method for determining Standard-Model inputs to improve our knowledge of nonperturbative meson gluon structure.
However, due to the notorious noise-to-signal issue associated with gluon observables, there have been only a few efforts to determine the first moment of pion gluon PDF~\cite{Meyer:2007tm,Shanahan:2018pib} with 450~MeV as the lightest pion mass calculated.
No kaon moments have yet been calculated.
There is little chance that we will be able to use lattice calculations of increasingly higher moments to reconstruct the $x$-dependence of the meson PDFs.
Recently, direct lattice calculations of the gluon PDF of hadrons have been proposed: pseudo-PDF~\cite{Balitsky:2019krf} and quasi-PDF~\cite{Zhang:2018diq,Wang:2019tgg} approaches, following the proposal of Large-Momentum Effective Theory (LaMET)~\cite{Ji:2013dva,Ji:2014gla,Ji:2017rah}.
Ongoing efforts using other methods, such as hadronic tensor~\cite{Liu:1993cv}, the operator product expansion~\cite{Detmold:2005gg}, lattice cross sections~\cite{Ma:2014jla,Ma:2017pxb} and the pseudo-PDF approach~\cite{Radyushkin:2017cyf,Orginos:2017kos}, have also had some success in extracting $x$-dependent hadron structure.
Most such PDF calculations, however, are done using the popular quasi-PDF and pseudo-PDF techniques and mostly limited to isovector quark distributions in the nucleon and the valence-quark distribution in the pion and kaons~\cite{Lin:2013yra,Lin:2014zya,Chen:2016utp,Lin:2017ani,Alexandrou:2015rja,Alexandrou:2016jqi,Alexandrou:2017huk,Chen:2017mzz,Alexandrou:2018pbm,Chen:2018xof,Chen:2018fwa,Alexandrou:2018eet,Lin:2018qky,Fan:2018dxu,Liu:2018hxv,Wang:2019tgg,Lin:2019ocg,Chen:2019lcm,Lin:2020reh,Chai:2020nxw,Bhattacharya:2020cen,Lin:2020ssv,Zhang:2020dkn,Li:2020xml,Fan:2020nzz,Gao:2020ito,Lin:2020fsj,Lin:2021brq,Zhang:2020rsx,Alexandrou:2020qtt,Alexandrou:2020zbe,Lin:2020rxa,Gao:2021hxl,Lin:2020rut,Orginos:2017kos,Karpie:2017bzm,Karpie:2018zaz,Karpie:2019eiq,Joo:2019jct,Joo:2019bzr,Radyushkin:2018cvn,Zhang:2018ggy,Izubuchi:2018srq,Joo:2020spy,Bhat:2020ktg,Fan:2020cpa,Sufian:2020wcv,Karthik:2021qwz,Chen:2018fwa,Sufian:2019bol,Izubuchi:2019lyk,Joo:2019bzr,Sufian:2020vzb,Shugert:2020tgq,Gao:2020ito}.

The first exploratory study applying the quasi-PDF approach to gluon PDFs~\cite{Fan:2018dxu} used overlap valence fermions on gauge ensembles with 2+1 flavors of domain-wall fermion and spacetime volume $24^3\times 64$, $a=0.1105(3)$~fm, and $M_\pi^\text{sea}=330$~MeV.
The gluon operators were calculated for all spacetime lattice sites and high statistics: 207,872 measurements were taken of the two-point functions with valence quarks at the light sea and strange masses (corresponding to pion masses 340 and 678~MeV, respectively).
Both pion and nucleon coordinate-space gluon quasi-PDF matrix-element ratios were compared with global fits after Fourier transformation;
the lattice results were consistent with global fits within large uncertainties.
Further progress was made using the pseudo-PDF method.
Fan et~al. studied the nucleon gluon PDF using clover fermions on $N_f=2+1+1$ MILC lattices, with a single lattice spacing 0.12~fm and two valence pion masses: 310 and 690~MeV~\cite{Fan:2020cpa}.
This was the first $x$-dependent gluon PDF calculated using lattice QCD.
Khan et~al. used distillation nucleon interpolating operators and extracted the nucleon gluon PDF on a 2+1-flavor 358-MeV pion mass ensemble at lattice spacing 0.094~fm~\cite{HadStruc:2021wmh}.
Both works neglected the mixing of the gluon operator with the quark-singlet sector.
Since then, Fan has continued to pursue gluon PDF calculations of the nucleon and pion using $N_f=2+1+1$ MILC lattices with multiple lattice spacings and pion masses.
Reference~\cite{Fan:2021bcr} reported the first lattice-QCD pion $x$-dependent PDFs at two lattice spacings: 0.15 and 0.12~fm and $M_\pi \approx 220$, 310 and 690~MeV);
the work used the global-fit PDF quark contribution to estimate the effect of the quark mixings in the gluon PDF determination.
However, kaon PDF calculations remain limited to the valence-quark distribution~\cite{Lin:2020ssv}.

In this work, we present the first lattice-QCD calculation of the $x$-dependent kaon gluon parton distribution using the pseudo-PDF method.
We study the kaon PDFs using two lattice spacings, 0.12 and 0.15~fm, and  one pion mass of 310~MeV.
The rest of the paper is organized as follows.
In Sec.~\ref{sec:lattice-details}, we discuss the operator used in this calculation, the numerical setup of lattice simulation, and how the ground-state kaon gluon matrix elements are obtained on the lattice.
In Sec.~\ref{sec:results}, we discuss our strategy for extracting the kaon gluon PDF from our lattice calculations.
The systematic uncertainties introduced by different steps are studied, and the lattice-spacing and pion-mass dependence are investigated.
